% !TeX root = ../main.tex
\chapter{Introduction}
\label{chap:intro}

\section{Context}
Computer scientists have always been seen as “hackers” that work with Command Line Interfaces (CLI). 
This belief comes from a long-gone era where Graphical User Interfaces (GUI) wouldn’t exist. In reality, 
Command Line Interfaces and Graphical User Interfaces live together in peace in the different sectors of informatics. 
However, lately the Graphical User Interfaces are experiencing, perhaps what can be called a ‘golden era’. 
The evolution from the first Graphical User Interface computer, the Xerox Alto in 1973\cite{wikipedia_gui},to Web GUIs applications 
still share the same core. However, they have quickly evolved in the past years thanks to eased access to the Internet 
and computers. Nowadays, big companies dictate the cutting edge of technology regarding GUIs. A trend that 
has become obvious is the rise of Dark mode, dividing the whole user community of many applications.

\section{Motivation}
Year after year we have seen how the whole world has changed at a really high pace. Two of the fastest 
changing areas are Graphical User Interfaces and User Experience (UX). A long time has passed since the 
times when developers would consider the front end a secondary part of development. 

Nowadays a new problem has become viral, not only inside the informatics world but throughout the whole 
spectrum of internet users: is dark mode better (from now on negative polarity) or bright mode (from now on 
positive polarity)? This has led many companies to embrace dark mode as its corporative colors, for example 
the social network X (former Twitter) shifted to negative polarity as its default mode. 

There have been many studies regarding usability between both options. However, no one has actually studied 
the effects of the positive and negative display polarity in longer texts. That has risen the following questions:
\begin{itemize}
    \item Does the user get more fatigued while reading in negative polarity?
    \item Will they read faster or slower?
    \item Will the user struggle to find more information?
    \item E-commerce websites are mainly programmed in positive polarity. Is the negative polarity really a disadvantage for the 
    user experience while reading and analyzing a product?
\end{itemize}

All these questions will be analyzed and given the best answer possible in this thesis. 
A conclusion regarding long texts and display polarities will be studied and explained here.

\section{Purpose of the thesis}
The purpose of the thesis is to investigate the visual fatigue and readability of both display polarities in terms of reading 
long texts in web-based products. This will be achieved by using a tobii eye-tracker. This tool can give real time data of the 
eye, giving us objective and solid information.

Visual fatigue can be defined as a moment in which the eyes have a huge workload, and the user can perceive a sensation of discomfort\cite{visualFatigue}. 
Visual fatigue can happen after a long period of time in front of a screen, but it also happens every time someone reads a text on a 
smaller scale. It is usually studied with questionnaires leading to subjectiveness in the analysis. However, in recent times, tools 
such as eye-trackers have proven a new way to study it. Symptoms of visual fatigue can be calculated as the time the user takes to read 
a whole paragraph and how many times the person takes to stop at each word.

\section{Document Structure}
The structure of the document has been briefly inspired by Jose Manuel Redondo’s spreadsheet of final thesis for development thesis 
\cite{plantilla_tfg}. However, in order to adapt it to an investigation thesis the spreadsheet has been adapted and now doesn't follow non specific standard.
This new “version” is mostly inspired by the scientific method by expanding the common topics on a scientific paper, while at the same time being 
mixed with the most common sections of a usual informatics thesis.

